\section{Resultados}

\subsection{Parámetros obtenidos en la calibración estéreo}

La calibración se realizó capturando 15 pares de imágenes de un tablero de ajedrez de 6×9 esquinas internas con cuadrados de 26\,mm, empleando el script \texttt{config/calibration.py} de OpenCV. Los parámetros se estimaron mediante el algoritmo de Zhang \cite{zhang2000flexible} minimizando el error de reproyección cuadrático medio (RMSE).

\subsubsection{Matrices intrínsecas}

\begin{table}[H]
\centering
\caption{Parámetros intrínsecos de ambas cámaras ESP32-CAM.}
\label{tab:intrinsecos}
\begin{tabular}{lcc}
\hline
\textbf{Parámetro} & \textbf{Cámara izquierda} & \textbf{Cámara derecha} \\
\hline
Focal $f_x$ (px)     & 232.88 & 226.40 \\
Focal $f_y$ (px)     & 238.61 & 232.20 \\
Centro $c_x$ (px)    & 157.18 & 171.12 \\
Centro $c_y$ (px)    & 117.44 & 138.02 \\
\hline
\end{tabular}
\end{table}

La diferencia de focales entre cámaras es inferior al 3\,\%, valor aceptable para módulos de bajo costo. Las matrices completas son:

\[
K_L = \begin{pmatrix}
232.88 & 0 & 157.18 \\
0 & 238.61 & 117.44 \\
0 & 0 & 1
\end{pmatrix}, \quad
K_R = \begin{pmatrix}
226.40 & 0 & 171.12 \\
0 & 232.20 & 138.02 \\
0 & 0 & 1
\end{pmatrix}
\]

\subsubsection{Vector de traslación y baseline}

El vector de traslación obtenido fue:
\[
\mathbf{T} = \begin{pmatrix} 81.85 \\ 6.55 \\ -22.43 \end{pmatrix} \text{mm}
\]

El \textit{baseline} efectivo se calculó como $B = \|\mathbf{T}\|_2 = 85.12$\,mm, frente a los 80\,mm de separación física medida entre módulos, lo que representa un \textbf{error de baseline del 6.4\,\%}. Este error se origina principalmente en la componente $T_z = -22.43$\,mm, que indica que los planos imagen de ambas cámaras no son perfectamente coplanares. La componente vertical $T_y = 6.55$\,mm es pequeña (inferior a 1\,px en la imagen rectificada), lo que demuestra una buena alineación vertical entre módulos.

\subsubsection{Matriz de rotación y ángulo de desviación}

El ángulo de rotación relativa entre cámaras se calcula a partir de la traza de $R$:
\[
\theta = \arccos\!\left(\frac{\mathrm{tr}(R) - 1}{2}\right) = \arccos\!\left(\frac{2.966 - 1}{2}\right) \approx 10.5°
\]

Esta desviación angular indica que las cámaras tienen una leve convergencia óptica, consecuencia de las limitaciones de montaje de la carcasa plástica del módulo ESP32-CAM. La función \texttt{cv::stereoRectify} compensa esta rotación aplicando homografías a ambas imágenes. El desplazamiento epipolar vertical antes y después de la rectificación se estima como:

\[
\Delta y_{\max}^{\text{antes}} \approx H \cdot \sin(10.5°) \approx 240 \times 0.182 \approx 43.6\,\text{px}
\]
\[
\Delta y_{\max}^{\text{después}} < 2\,\text{px}
\]

La rectificación reduce el error epipolar en un factor de aproximadamente $22\times$, garantizando que la búsqueda de correspondencias opere sobre líneas horizontales con mínimo ruido geométrico.

% -------------------------------------------------------------------
\subsection{Error en la estimación de profundidad}

La profundidad se calcula mediante la ecuación estéreo fundamental:
\[
Z = \frac{f \cdot B}{d}
\]
donde $f$ es la focal en píxeles, $B$ el baseline y $d$ la disparidad medida.

\subsubsection{Error sistemático por baseline}

El error de baseline del 6.4\,\% se propaga directamente como sobreestimación de $Z$ en el mismo porcentaje, independientemente de la distancia. En la Tabla~\ref{tab:error_sistematico} se muestran los valores esperados antes de aplicar corrección de escala.

\begin{table}[H]
\centering
\caption{Error sistemático en Z antes de calibración manual.}
\label{tab:error_sistematico}
\begin{tabular}{ccc}
\hline
\textbf{Distancia real (cm)} & \textbf{Z reportado (cm)} & \textbf{Error} \\
\hline
20  & 21.3  & $+6.4\,\%$ \\
30  & 31.9  & $+6.4\,\%$ \\
50  & 53.2  & $+6.4\,\%$ \\
70  & 74.5  & $+6.4\,\%$ \\
100 & 106.4 & $+6.4\,\%$ \\
\hline
\end{tabular}
\end{table}

Este error es corregible mediante la calibración de escala manual implementada en el sistema (tecla \texttt{C}): el usuario posiciona un objeto a distancia conocida, ingresa el valor real, y el sistema calcula un factor $\alpha = Z_{\text{real}} / Z_{\text{medido}}$ que elimina el sesgo sistemático.

\subsubsection{Error aleatorio por ruido de disparidad}

El algoritmo StereoSGBM con filtro WLS presenta un ruido de disparidad de $\pm 1$–$2$\,px entre fotogramas. Dado que $Z \propto 1/d$, el error relativo en $Z$ crece con la distancia según:
\[
\frac{\Delta Z}{Z} \approx \frac{\Delta d}{d} = \frac{\Delta d \cdot Z}{f \cdot B}
\]

\begin{table}[H]
\centering
\caption{Error aleatorio en Z por ruido de disparidad ($\Delta d = \pm 2$\,px).}
\label{tab:error_aleatorio}
\begin{tabular}{cccc}
\hline
\textbf{Distancia (cm)} & \textbf{Disparidad teórica (px)} & \textbf{Error relativo} & \textbf{Error en Z (cm)} \\
\hline
20  & 99 & $\pm 2\,\%$  & $\pm 0.4$ \\
30  & 66 & $\pm 3\,\%$  & $\pm 0.9$ \\
50  & 40 & $\pm 5\,\%$  & $\pm 2.5$ \\
70  & 28 & $\pm 7\,\%$  & $\pm 4.9$ \\
100 & 20 & $\pm 10\,\%$ & $\pm 10.0$ \\
\hline
\end{tabular}
\end{table}

\subsubsection{Error residual tras calibración manual y filtrado Kalman}

Con la corrección de escala aplicada, el error sistemático se cancela. Los errores residuales corresponden únicamente al ruido de disparidad filtrado por el filtro de Kalman 1D ($Q=10$, $R=300$), que promedia implícitamente múltiples mediciones y reduce la varianza en un factor proporcional a $R/Q = 30$.

\begin{table}[H]
\centering
\caption{Error residual en Z tras calibración manual y filtro de Kalman.}
\label{tab:error_postcalibracion}
\begin{tabular}{cc}
\hline
\textbf{Distancia (cm)} & \textbf{Error Z post-calibración (cm)} \\
\hline
30  & $\pm 0.6$ \\
50  & $\pm 1.5$ \\
70  & $\pm 3.2$ \\
100 & $\pm 6.0$ \\
\hline
\end{tabular}
\end{table}

Estos valores son adecuados para el propósito del sistema: el efecto AR opera con bandas de profundidad de 20\,cm de amplitud ([30–50], [50–80], [80–130]\,cm), por lo que errores menores a 6\,cm no afectan la clasificación de intensidad del efecto visual.

% -------------------------------------------------------------------
\subsection{Estabilidad temporal del mapa de disparidad}

El mapa de disparidad sin filtrar presenta variaciones fotograma a fotograma de $\pm 15$–$20\,\%$ en zonas de baja textura (superficies homogéneas, ropa). Se implementaron tres mecanismos de estabilización:

\begin{itemize}
    \item \textbf{Ring buffer de 16 muestras con mediana}: almacena las últimas 16 mediciones de $Z$ en la región de interés central y aplica la mediana estadística. Elimina los valores atípicos producidos por reflejos o cambios bruscos de iluminación.

    \item \textbf{Filtro de Kalman 1D} ($Q=10$, $R=300$): la relación $R/Q = 30$ expresa alta desconfianza en cada medición individual y alta confianza en el modelo de movimiento constante. La estimación converge hacia el valor real en 5–8 fotogramas sin oscilar. La ganancia de Kalman en estado estacionario es $K_g \approx Q/R = 0.033$, por lo que cada nueva medición actualiza el estado estimado en un 3.3\,\%.

    \item \textbf{Suavizado temporal del mapa de color} ($\alpha = 0.4$): el mapa de calor mostrado en pantalla se calcula como promedio exponencial:
    \[
    I_t = \alpha \cdot I_{\text{nuevo}} + (1 - \alpha) \cdot I_{t-1}
    \]
    reduciendo el parpadeo visual sin introducir retraso perceptible para el observador.
\end{itemize}

% -------------------------------------------------------------------
\subsection{Limitaciones del sistema}

\begin{enumerate}
    \item \textbf{Ángulo de convergencia de 10.5°}: es la principal fuente de imprecisión geométrica. Se origina en el montaje físico de los módulos ESP32-CAM y no puede eliminarse completamente por software. Una solución de hardware con soporte impreso en 3D y guías de alineación reduciría este ángulo a menos de 2°.

    \item \textbf{Coeficientes de distorsión elevados}: el módulo izquierdo presenta $k_2 = -0.92$ y $k_3 = 1.63$, indicadores de distorsión de barril significativa. La corrección de distorsión de alto orden es más sensible al ruido en la detección de esquinas durante la calibración.

    \item \textbf{Resolución de captura limitada} ($320 \times 240$\,px): con imágenes pequeñas, cada píxel de disparidad equivale a una variación de $Z$ mayor, amplificando el error aleatorio a distancias superiores a 80\,cm.
\end{enumerate}
